\chapter{Constraint Projection}

\section{Where Proposals Go to Be Checked}

Return to the proposal Chapter 12 used to illustrate what addition alone cannot catch: a \(\Delta\psi\) that asserts the courtyard's ground is simultaneously undisturbed and freshly dug, fluent enough to have been generated by a well-trained \(F_\psi\), selected with high relevance by the policy Chapter 13 described, and still, on its own terms, incoherent. Nothing in the last two chapters explains what stops this proposal from becoming part of the field. This chapter explains that.

\section{The Constraint Manifold}

Every rule this book has stated since Chapter 8 has, without using the word, been describing a constraint: a condition a field state must satisfy to be admissible. It is worth collecting them now under a single name. Write C for the constraint manifold, the set of field states this book's framework regards as admissible. C is not introduced here for the first time; it has been accumulating members since Part II began, and this section's task is only to name what was already implicit.

Chapter 8 contributed the entropy constraint: no admissible state may show a local entropy decrease at a point outside the region actually observed during the update that produced it. Chapter 9 contributed the invention-to-memory discipline: no admissible state may hold an appearance at a settled point that discards what was previously committed there without cause. Chapter 11 contributed residue consistency: no admissible state may hold an appearance or entropy value that contradicts what the point's own accumulated residue has already ruled out, undisturbed paving where the residue records recent digging. Chapter 10's bridge between cue coherence and field coherence, not yet formalized but already promised, will eventually contribute a further member, ruling out states where an agent's realized action fails to correspond to any admissible field transformation at all.

C, in other words, is not a single equation waiting to be written down in this chapter. It is the intersection of every discipline this book has argued for since Chapter 8, growing as the book's argument grows, and this chapter's job is not to complete it, which properly belongs to Chapter 18's fuller formal treatment, but to explain what happens once a proposal is checked against whatever of it has already been established.

\section{Projection as Nearest Admissible State}

Given a proposal \(M_t\) + \(\Delta\psi\) that may or may not lie in C, projection does not simply accept or reject it. It finds the nearest point in C to what was proposed:

\[ Y = \operatorname*{argmin}_{Y \in C} \|Y - (M_t + \Delta\psi)\|^2. \]

Read this as a kind of semantic pressure solve, in the sense a physical simulation enforces incompressibility or a mechanical system enforces a joint constraint: rather than discarding a proposal that oversteps what is admissible, the system finds the closest admissible configuration to the one that was actually suggested, preserving as much of the proposal's intent as the constraints allow. Write \(\Pi_C\) for this operation as a whole, so that Y = \(\Pi_C\)(\(M_t\) + \(\Delta\psi\)). The contradictory courtyard proposal from section 14.1 does not simply vanish under this operation. It is pulled toward whatever admissible resolution lies nearest to it, most plausibly the resolution the residue actually supports, freshly dug ground with an appearance to match, discarding only the part of the proposal that residue had already ruled out, and keeping whatever else the proposal specified that did not conflict.

This is worth dwelling on because it explains something about the character of this framework that could otherwise seem harsh. Constraint projection is not a censor that vetoes proposals outright. It is closer to a correction: it trusts the proposal's direction and intent, and adjusts only what admissibility requires adjusting. A neural proposal that gets almost everything right and stumbles on one incompatible detail will emerge from projection nearly intact, missing only the detail that could not survive contact with what the field already holds.

\section{Revisiting the Two Sources}

Chapter 13 distinguished \(\Delta\psi\)'s two sources, the field's intrinsic relaxation \(F_{\mathrm{phys}}\) and the learned neural forcing \(F_\psi\), and suggested, without yet justifying the claim, that the first is close to admissible almost by construction while the second carries most of a proposal's risk. Constraint projection is where that claim can finally be earned rather than merely asserted.

\(F_{\mathrm{phys}}\) is defined as the field's own rule-governed tendency to change, the same tendency that produces entropy relaxation, coherence smoothing, and the ordinary drift of residue over time. Because these tendencies are themselves derived from the same discipline that defines C, an entropy field that relaxes according to its own dynamics does not thereby violate the entropy constraint, and a coherence field that smooths according to its own dynamics does not thereby violate coherence consistency. \(F_{\mathrm{phys}}\), in other words, is constructed to already respect C, in the same sense that a ball rolling under gravity along a track built to constrain it does not need a separate check to confirm it stays on the track. Projected, \(F_{\mathrm{phys}}\) typically maps close to itself, and the distance ‖Y − (\(M_t\) + \(\Delta\psi\))‖ contributed by this term alone is usually small.

\(F_\psi\) carries no such guarantee, and this is the honest version of the claim Chapter 13 made informally. A neural forcing term is trained to be expressive, to propose changes that a hand-specified \(F_{\mathrm{phys}}\) never could, sitting on a bench, moving through a doorway, inventing an appearance for a previously unobserved patch of ground, and expressivity of this kind is exactly what makes a proposal likely to reach outside C. Most of the actual work constraint projection performs, on a typical update, is work performed on the \(F_\psi\) contribution to \(\Delta\psi\), correcting or trimming what the learned, expressive half of the proposal got wrong, while the intrinsic half passes through largely unchanged. The asymmetry Chapter 13 flagged was correct, and this section is where it stops being an informal impression and becomes a consequence of how \(F_{\mathrm{phys}}\) and C were built from the same discipline in the first place.

\section{What Projection Cannot Fix}

Nearest-point projection assumes that some admissible state is worth moving toward, that C contains something in the vicinity of what was proposed. This assumption can fail. A proposal sufficiently incompatible with everything the field currently holds, not merely a locally contradictory detail but a wholesale departure from the courtyard's accumulated coherence, entropy, and residue all at once, may have no nearby admissible state worth settling for; the nearest point in C might be so distant from the original proposal that accepting it would honor almost none of what the proposal actually intended. Projection, as this chapter has developed it, still returns some answer in this case, the mathematics of the minimization does not fail. But an answer this distant from what was proposed is arguably not a correction of the original intent at all, and treating it as one risks quietly replacing a bad proposal with an unrelated one under the appearance of continuity. This book will not resolve that tension here. It belongs to the same family of questions as outright refusal, the deliberate rejection of a proposal rather than its correction, and refusal has been mentioned only in passing, in the Spherepop-derived material this book draws on elsewhere, as an irreversible elimination of branches rather than a repair performed on one. Whether and when projection should give way to refusal is a question this book sets aside for now and returns to only once Part IV has given constraint projection its full mathematical treatment.

\section{Looking Ahead}

This chapter has given proposals somewhere to go once they have been generated: projected onto the nearest state the field's accumulated constraints actually permit, with the field's own intrinsic tendencies passing through nearly untouched and a proposal's learned, expressive content bearing most of the correction. What this chapter has not yet done is explain how a projected state, Y, actually becomes \(M_{t+1}\). The two are not automatically the same thing, and Chapter 15 is where that distinction, promised as far back as this book's roadmap, is finally made precise.
